\chapter{Development Plan}
\ChapterLabel{DevelopmentPlan}
\section{Chapter Responsibility}
\label{Section:DevelopmentPlanResponsibility}
Define how the project is governed, planned, reviewed, and changed.

\noindent\textbf{Expected result:} An actionable plan with visible ownership and decision gates.

\section{Project Governance}
\label{Section:ProjectGovernance}
Joel Murphy and Jason Lefler own this document, approve revisions and the initial architecture
and context baseline, and resolve all open \OwnedTBD{} items. The owner completed the
gap-list review and authorized these document updates. That authorization does
not resolve the retained \OwnedTBD{} items, accept residual risk, establish a detailed
requirements baseline, or authorize release. Additional acceptance criteria
are \OwnedTBD{}. All open risks also require owner assignment and assessment.

\subsection{Ownership of Unresolved Decisions}
\label{Section:TBDDecisionOwnership}
Every unresolved decision is marked TBD or \OwnedTBD{}. Joel Murphy and Jason Lefler are
the named decision co-owners for every such marker, including those without an
inline owner name. This ownership identifies
the people responsible for resolving the decision; it does not assign an
undecided hardware, firmware, host-integration, toolchain, maintenance, or risk
role to either person. Joel Murphy and Jason Lefler are the named decision co-owners for all open
items, including items in prose, tables, metadata, and retained template
guidance. Items remain unresolved until an explicit decision is recorded.

\subsection{Requirements Completion Gate}
\label{Section:RequirementsCompletionGate}
The original two story-derived proposals are retained and supplemented by
candidates translating the owner-approved acquisition and interruption
decisions. The owner authorized these source edits without compilation or
Git operations; implementation, baseline approval and verification remain pending. \hyperref[Table:ProjectProfileConcernCoverage]{The embedded
concern coverage register} identifies mapped, partially mapped, and open
concerns and the previously authorized Not Applicable dispositions.
AIS-06 is not complete while any applicable concern lacks a requirement
mapping or an approved Not Applicable rationale. Joel Murphy and Jason Lefler own every open
decision. Producing a document PDF does not close this gate.

\section{Domain Profile Selection}
\label{Section:DevelopmentPlanProfiles}
The available profile was read before editing, and its concerns were added to
the reviewed gap list. The project facts remain in their core chapters.
\small
\begin{tabularx}{\linewidth}{L{3.0cm}L{2.2cm}Y}
\caption{Project Domain Profile Decisions\label{Table:ProjectDomainProfiles}}\\
\toprule
\textbf{Profile} & \textbf{Status} & \textbf{Rationale, Owner, and Date}\\
\midrule
\endfirsthead
\multicolumn{3}{l}{\small\itshape Project Domain Profile Decisions (continued)}\\
\toprule
\textbf{Profile} & \textbf{Status} & \textbf{Rationale, Owner, and Date}\\
\midrule
\endhead
\bottomrule
\endlastfoot
\hyperref[Chapter:EmbeddedSystemsProfile]{Embedded Systems} & Applicable & Confirmed by Joel Murphy (co-owner from Version 0.2: Jason Lefler). Applicability triggers are documented in \hyperref[Section:EmbeddedApplicabilityTriggers]{Embedded Applicability Triggers}. Recorded in revision 0.1, November 8, 2026.\\
\end{tabularx}
\normalsize

Only the Embedded Systems profile is implemented in the supplied template.
The profile adds hardware and toolchain ownership, timing and resource review,
controlled hardware evidence, and target-test review concerns. Detailed gates,
test roles, and evidence criteria are deferred with their dependent artifacts.
Future requirements work must map each applicable profile concern to a
requirement or an approved Not Applicable disposition under
\hyperref[Section:DomainProfileRequirements]{the domain-profile translation rule}.

\section{Embedded Applicability Triggers}
\label{Section:EmbeddedApplicabilityTriggers}
The owner identified the following triggers as the justification for applying
\hyperref[Chapter:EmbeddedSystemsProfile]{the Embedded Systems profile} to
the sensor-interface subsystem, meaning the FTDI245R acquisition pipeline
code base. The recorded decision remains Applicable.
\begin{itemize}
\item \textbf{Dedicated interface hardware:} the acquisition context depends
on external hardware, including the
\hyperref[Dependency:ADC0804]{ADC0804},
\hyperref[Dependency:Counter]{74xx counter}, and
\hyperref[Dependency:PeerFTDI245R]{peer FTDI245R}.
\item \textbf{Sensors and real-time events:} the Sensor and acquisition events
make event handling a relevant embedded concern. Event definitions and any
real-time guarantees remain \OwnedTBD{}.
\item \textbf{Host-OS driver or library interaction:} the pipeline depends on
the \hyperref[Dependency:USBStack]{USB stack} and
\hyperref[Dependency:HostOS]{host operating system}; the driver/API and library
version and confirmed binding remain TBD (decision co-owners: Joel Murphy and Jason Lefler);
FTD2XX is proposed and the D2XX Guide v1.3 is the documentation baseline.
\item \textbf{Sampling and conversion timing:} timing at the Sensor,
ADC0804, and acquisition interface requires review. Sample rates, conversion
timing, triggering, latency, and jitter limits remain \OwnedTBD{}.
\item \textbf{Bounded host resources:} processing, buffering, memory,
storage, and transfer capacity require resource review. Allocations, limits,
margins, and supporting measurements remain \OwnedTBD{}.
\end{itemize}
These triggers establish the need for tailoring; they do not establish
hardware compatibility, timing guarantees, resource sufficiency, or test results.

\section{Hardware and Firmware Responsibilities}
\label{Section:HardwareFirmwareResponsibilities}
Joel Murphy and Jason Lefler retain human approval authority and resolve the \OwnedTBD{} decisions
below. Decision ownership does not assign hardware or firmware engineering
responsibilities to him. Technical owners and implementation allocation have
not yet been provided.
\small
\begin{tabularx}{\linewidth}{L{3.0cm}L{2.2cm}Y}
\caption{Hardware and Firmware Responsibility Status\label{Table:HardwareFirmwareResponsibilities}}\\
\toprule
\textbf{Area} & \textbf{Technical Owner} & \textbf{Responsibility / Allocation Status}\\
\midrule
\endfirsthead
\multicolumn{3}{l}{\small\itshape Hardware and Firmware Responsibility Status (continued)}\\
\toprule
\textbf{Area} & \textbf{Technical Owner} & \textbf{Responsibility / Allocation Status}\\
\midrule
\endhead
\bottomrule
\endlastfoot
Hardware & \OwnedTBD{} & Assignment of responsibility for hardware identification,
revisions, wiring, electrical compatibility, sampling/conversion timing,
and controlled hardware evidence is \OwnedTBD{}. Hardware remains external under
\hyperref[Section:ProjectSystemBoundary]{the system boundary}.\\
Firmware & \OwnedTBD{} & Whether firmware is required, where it executes, and who
owns its implementation, configuration, and evidence are \OwnedTBD{}. No firmware
component, platform, or engineering assignment is established.\\
Host integration & \OwnedTBD{} & Assignment of responsibility for driver/library
selection and Linux, Windows, and macOS integration and compatibility evidence
is \OwnedTBD{}. The \hyperref[Table:ContextDependencies]{dependency register} defines
the controlled external context.\\
\end{tabularx}
\normalsize

\section{Toolchain and Host-Platform Assumptions}
\label{Section:ToolchainHostAssumptions}
Linux, Windows, and macOS are the owner-confirmed host platforms in the
initial context baseline. Supported operating-system versions, architectures,
permissions, USB-stack implementations, actual FTD2XX driver/library version,
and compatibility evidence remain TBD (decision co-owners: Joel Murphy and Jason Lefler). FTD2XX
is the proposed library; D2XX Programmer's Guide v1.3 is the approved
documentation baseline, not an installed software version. No successful cross-platform execution
is claimed.

Programming languages, development tools, compiler/linker versions, SDKs,
dependency versions, build environment, and reproducibility acceptance
criteria remain \OwnedTBD{}. Toolchain ownership is \OwnedTBD{}. Git and GitHub configuration
control is recorded in
\hyperref[Section:ProjectConfigurationControl]{Configuration Control and Service Dispositions};
it does not establish a selected implementation toolchain. No bare-metal,
RTOS, embedded Linux, programmable-logic, or firmware execution platform is
assumed. Detailed allocation remains governed by
\hyperref[Section:ProjectSystemBoundary]{System Boundary and Allocation}.

\section{Configuration Control and Service Dispositions}
\label{Section:ProjectConfigurationControl}
Configuration control uses Git, with GitHub for remote repository hosting.
The repository location, controlled source revision, build environment, and
reproducibility acceptance criteria are \OwnedTBD{}. Git and GitHub alone do not
demonstrate that a code base is reproducible.

\subsection{Sustainment Notes}
\label{Section:ProjectSustainmentNotes}
Joel Murphy (co-owner from Version 0.2: Jason Lefler) marked updates, rollback, calibration records, manufacturing,
field service, and component availability Not Applicable because this system
will not be field deployed. These are recorded owner dispositions; a scope
change may require their review. This disposition does not establish a
technical claim about whether calibration is needed for a particular sensor.

\section{Development Sequence}
\label{Section:ProjectDevelopmentSequence}
This revision contains context, an initial capability baseline, and reviewed
gaps. Later work proceeds from context and stakeholder needs through detailed
requirements, use cases, architecture views, tests and traceability, evidence,
and release. Resolve consequential hardware, boundary, interface, and
acceptance gaps before dependent detailed work. Week assignments and dates
beyond the current revision are \OwnedTBD{}. No implementation or test pass is claimed.

\section{Reviewed Gap List}
\label{Section:ProjectGapReview}
The owner reviewed G01--G16 one item at a time. Review is complete; the open
items below remain unresolved. Joel Murphy and Jason Lefler are the decision co-owners for every
\OwnedTBD{} item. Where no decision date was supplied, November 8, 2026 identifies the
revision record date rather than an independently evidenced approval timestamp.
\small
\begin{tabularx}{\linewidth}{L{1.1cm}L{3.0cm}Y}
\caption{Reviewed Project Gap List\label{Table:ProjectGapReview}}\\
\toprule
\textbf{ID} & \textbf{Area} & \textbf{Reviewed Disposition / Remaining Gap}\\
\midrule
\endfirsthead
\multicolumn{3}{l}{\small\itshape Reviewed Project Gap List (continued)}\\
\toprule
\textbf{ID} & \textbf{Area} & \textbf{Reviewed Disposition / Remaining Gap}\\
\midrule
\endhead
\bottomrule
\endlastfoot
\phantomsection\label{Gap:G01}G01 & Purpose & Project purpose and mission need recorded. Owner-approved \gls{fsareliable} and \gls{fsaplugandplay} definitions apply. Quantitative reliability targets remain TBD (decision co-owners: Joel Murphy and Jason Lefler).\\
\phantomsection\label{Gap:G02}G02 & Scope and boundary & Revision scope recorded in \hyperref[Section:ProjectScope]{Scope}. The pipeline code-base boundary and controlled external dependencies are defined in \hyperref[Section:ProjectSystemBoundary]{System Boundary and Allocation}. Detailed allocations remain \OwnedTBD{}.\\
\phantomsection\label{Gap:G03}G03 & Stakeholders & Owner, users, and approver recorded in \hyperref[Section:ProjectStakeholders]{Audience}. Additional assignments and organization are \OwnedTBD{}. Hardware, firmware, and host-integration responsibility gaps are recorded in \hyperref[Table:HardwareFirmwareResponsibilities]{Responsibility Status}.\\
\phantomsection\label{Gap:G04}G04 & Hardware baseline & Specific FTDI245R module and board revision are \OwnedTBD{}. Platform class, other hardware revisions, and controlled hardware evidence remain \OwnedTBD{}.\\
\phantomsection\label{Gap:G05}G05 & Signal definition & Sensor through ADC0804; one unsigned 8-bit byte per conversion, no headers/framing; constant sampling without source-rate bursts. Sample rate is TBD pending ADC0804 timing/hardware confirmation. Signal type, channels and triggering remain TBD (decision co-owners for each: Joel Murphy and Jason Lefler).\\
\phantomsection\label{Gap:G06}G06 & Interfaces & Connection chain and Linux/Windows/macOS retained. Raw-byte delivery occurs at consumer interface acceptance. FTD2XX proposed; D2XX Guide v1.3 is the documentation baseline. Actual library/driver version, OS versions, electrical/handshake contracts and concrete consumer API remain TBD (decision co-owners: Joel Murphy and Jason Lefler).\\
\phantomsection\label{Gap:G07}G07 & Constraints and toolchain & Programming languages, development tools, budget, schedule, licensing rules, and applicable standards are \OwnedTBD{}. See \hyperref[Section:ToolchainHostAssumptions]{Toolchain and Host-Platform Assumptions}.\\
\phantomsection\label{Gap:G08}G08 & Quality attributes & Continuous acquisition until user stop; \gls{fsareliable} and \gls{fsacorruption} definitions approved. Sample rate, acceptance-test duration, quantitative reliability, latency and accuracy targets remain TBD (decision co-owners: Joel Murphy and Jason Lefler).\\
\phantomsection\label{Gap:G09}G09 & Resources, power, environment & Buffering, storage, bandwidth, processing, power, temperature, and operating-condition limits are \OwnedTBD{}.\\
\phantomsection\label{Gap:G10}G10 & Lifecycle behavior & Automatic detection/start after installation/configuration; positive reads; partial/empty timeout continuation; host-read byte preservation, interruption report and automatic reconnect/resume approved. Device-resident/disconnected-source bytes excluded. Timeout duration, reporting API, stop/drain and detailed recovery contracts remain TBD (decision co-owners: Joel Murphy and Jason Lefler).\\
\phantomsection\label{Gap:G11}G11 & Safety and security & Hazards, electrical protection, safe states, access control, sensitive-data handling, and assurance dispositions are \OwnedTBD{}.\\
\phantomsection\label{Gap:G12}G12 & Updates and service & Git/GitHub and owner Not Applicable dispositions are recorded in \hyperref[Section:ProjectConfigurationControl]{Configuration Control and Service Dispositions}.\\
\phantomsection\label{Gap:G13}G13 & Risks & Six open risks are recorded in \hyperref[Table:ProjectOpenRisks]{Open Risks}. Assessment, mitigation, and risk ownership remain \OwnedTBD{}.\\
\phantomsection\label{Gap:G14}G14 & Evidence & FTDI TN\_167 v1.0 and D2XX Guide v1.3 plus owner decisions support candidate obligations. Software and physical FT245R sequence paths are planned, not executed. Hardware evidence, controlled baselines and objective results remain TBD (decision co-owners: Joel Murphy and Jason Lefler).\\
\phantomsection\label{Gap:G15}G15 & Approvals & Joel Murphy and Jason Lefler approve revisions and the initial baseline and resolve all open \OwnedTBD{} items. Additional acceptance criteria remain \OwnedTBD{}; see \hyperref[Section:ProjectGovernance]{Project Governance}.\\
\phantomsection\label{Gap:G16}G16 & Document setup & Namespace fsa; co-owners Joel Murphy and Jason Lefler; organization \OwnedTBD{}; revision 0.2; Draft --- Initial Architecture and Context Baseline; November 8, 2026. PDF compilation is assigned to the owner.\\
\end{tabularx}
\normalsize

\section{Open Risks}
\label{Section:ProjectOpenRisks}
The following candidate risks were accepted for listing as open risks during
the gap review. Their listing does not establish occurrence, severity,
likelihood, mitigation effectiveness, or residual-risk acceptance. Assessment,
mitigation, and risk ownership are \OwnedTBD{} for each risk.
\small
\begin{tabularx}{\linewidth}{L{1.4cm}L{4.1cm}Y}
\caption{Open Project Risks\label{Table:ProjectOpenRisks}}\\
\toprule
\textbf{ID} & \textbf{Open Risk} & \textbf{Concern to Assess}\\
\midrule
\endfirsthead
\multicolumn{3}{l}{\small\itshape Open Project Risks (continued)}\\
\toprule
\textbf{ID} & \textbf{Open Risk} & \textbf{Concern to Assess}\\
\midrule
\endhead
\bottomrule
\endlastfoot
\phantomsection\label{Risk:RSK001}RSK-001 & Throughput shortfall & Acquisition demand may exceed sustained pipeline capacity.\\
\phantomsection\label{Risk:RSK002}RSK-002 & Sample loss & Samples may be lost or overwritten without an agreed loss policy.\\
\phantomsection\label{Risk:RSK003}RSK-003 & Timing ambiguity & Unspecified sampling, timing, or triggering may make acquired data ambiguous.\\
\phantomsection\label{Risk:RSK004}RSK-004 & Electrical incompatibility & Unconfirmed connections and electrical limits may prevent correct or safe interaction.\\
\phantomsection\label{Risk:RSK005}RSK-005 & Host interruption & Disconnection or interrupted host activity may disrupt acquisition.\\
\phantomsection\label{Risk:RSK006}RSK-006 & Unsupported hardware assumptions & Uncontrolled hardware assumptions may invalidate design and test conclusions.\\
\end{tabularx}
\normalsize

\section{Execution and Review Record}
\label{Section:ProjectExecutionRecord}
This initialization applies \hyperref[Skill:AIS-02]{AIS-02} through
\ChapterRef{AIImplementationWorkflow}, with planning under
\hyperref[Skill:AIS-01]{AIS-01}. One agent acts as aiaCoordinator and
aiaDeveloper, followed by self-review. No subagents or independent review are
claimed. Sequential execution avoids coordination overhead for dependent
changes to shared chapters and registers.

Planning ranges for the approved update were: metadata and context,
20--40 minutes / 3,000--6,000 tokens; reviewed gaps, profile, risks, and
provenance, 20--40 minutes / 3,000--6,000 tokens; source consistency review,
10--20 minutes / 1,000--3,000 tokens. These are rough planning ranges, not
measured usage or cost guarantees. The owner elected to compile the PDF.
Compilation, build-log review, PDF link checks, and visual inspection remain
pending owner execution; source checks alone do not satisfy build acceptance.
